\section{Model Evaluation} \label{sec:model-evaluation}
After training, model performance was evaluated on the test set clips to ensure that seizure prediction was substantially better than chance. 
Only under this condition do the \glspl{fp} provide meaningful information about the model’s seizure propensity; \glspl{fp} from a near‑random classifier would not.


The models were employed to assign scores in the range [0, 1] to all existing segments, representing the model-derived probability that a given segment is preictal.  
Subsequently, the segment-level scores were averaged within each clip to obtain a corresponding clip-level score.  
These clip-level scores, rather than the raw segment-level scores, were used in subsequent evaluations.
Two types of metrics were then used to evaluate the performance: clip-based metrics and \glsreset{ebm}\glspl{ebm}.

Here, clip-based metrics refer to the standard measures used to evaluate models in machine learning on samples (clips), which use the confusion matrix of \glspl{tp}, \glsreset{fp}\glspl{fp}, \glspl{fn}, and \glspl{tn} to compute metrics, such as precision, sensitivity, accuracy, and the \gls{roc auc}.
The Hanley McNeil test \parencite{hanley_method_1983} was performed to test for statistically significant improvement of the \gls{roc auc} compared to a random classifier.

On the other hand, \glspl{ebm} are specific to time-series data and seizure prediction and quantify more clinically relevant properties.
Here, the (relative) \gls{tifw} \parencite{mormann_seizure_2007, eberlein_machine_2025} and the (relative) number of seizures predicted were used.
They are discussed in the following sections.

\subsection{Time in False Warning and Event-Based Specificity} \label{subsec:tifw}
When a system unnecessarily alarms the patient, it causes them distress and inconvenience.
It is therefore essential to minimize false alarms and the derived \gls{tifw}.

There is a caveat to calculating the \gls{tifw}: when a prediction is issued based on a clip, the seizure is predicted to occur during that clip's \gls{sph} ($\sim$\num{5}--\SI{65}{\minute} after the clip's end), rather than during the clip.
Since clip duration ($\sim$\SI{10}{\minute}) is not equal to \gls{sph} duration, summing the duration of \gls{fp} clips is not sufficient to calculate the \gls{tifw}.
The clip's \glspl{sph} and overlaps thereof must be taken into account.

Thus, the \gls{tifw} was calculated as follows:
\begin{enumerate}
    \item Only \textbf{valid, non-preictal clips} were selected.
    \item For all such clips that issued a warning (\glspl{fp}), the \textbf{\glspl{sph}} were selected.
    \item \textbf{Overlapping \glspl{sph}} were then merged, so that shared time was not counted multiple times.
    \item The total duration of these non-overlapping false-warning intervals was the \textbf{absolute \gls{tifw}}.
    \item The \textbf{relative \gls{tifw}} was obtained by normalizing absolute \gls{tifw} by the maximum possible warning time, which was the sum of the non-overlapping \gls{sph} intervals from \textbf{all} non-preictal clips.
\end{enumerate}
Relative \gls{tifw} is the event-based equivalent to the \textbf{\gls{fpr}}, i.e., the inverse of \textbf{specificity}.
That is to say, it is minimized when non-preictal \textbf{(false) samples} are \textbf{correctly classified} as such.
It is a metric bounded in [0, 1], enabling comparable interpretation across patients and consistent with the guidelines in \textcite{mormann_seizure_2007}.


\subsection{Event-Based Sensitivity} \label{subsec:eb-sensitivity}
It is also essential to quantify the \textbf{relative number of seizures successfully predicted}, which is also referred to as \glsreset{eb}\gls{eb} sensitivity.
For preictal clips, a positive prediction is considered a true warning, and a seizure is counted as predicted if it falls within the \gls{sph} of at least one positively predicted preictal clip.
This yields the number of predicted seizures, and \gls{eb} sensitivity is then computed as the ratio of predicted to total seizures.
This metric therefore reflects clinically relevant seizure prediction performance at event-based level rather than clip level.

\subsection{Threshold Optimization} \label{subsec:threshold-optimization}
Whether a clip issued a warning depended on whether its score was greater than or equal to a classification threshold.
To optimize it, metrics were calculated for all possible thresholds, i.e., all unique model scores.
To take both the \gls{eb} specificity (1 - rel. \gls{tifw}) and the \gls{eb} sensitivity into account, the harmonic mean $H$ between them was calculated:
\begin{equation} \label{eq:harmonic-mean}
    H(a,b)=\frac{2ab}{a+b}
\end{equation}
\begin{equation} \label{eq:eb-score}
    \mathrm{Score}_{\mathrm{EB}} = H\!\left(\mathrm{Sensitivity}_{\mathrm{EB}},\, 1-\mathrm{rel.\ TIFW}\right)
\end{equation}
The threshold that \textbf{maximized this \gls{eb} score} was set as the \textbf{best threshold} and used to determine binary clip predictions in further steps.

The metrics were calculated independently on the training and test sets.
In causal, \mbox{(pseudo-)prospective} study designs, classification thresholds should be determined exclusively using the training set.
However, as this work prioritizes the cycle-based evaluation over a realistic measure of performance, the threshold was individually optimized for the training and test sets, which boosted performance on the test set.
